\section{Results}
\label{sec:results}

\begin{table*}[t]
  \centering\small\setlength{\tabcolsep}{4pt}
  \begin{tabular}{@{}lcccccccc@{}}
\toprule
 & \multicolumn{2}{c}{RESISC45} & \multicolumn{2}{c}{EuroSAT-MS} & \multicolumn{2}{c}{PlantVillage} & \multicolumn{2}{c}{PlantDoc} \\
\cmidrule(lr){2-3}\cmidrule(lr){4-5}\cmidrule(lr){6-7}\cmidrule(lr){8-9}
Score at 30\% of the parameters & P & PMP & P & PMP & P & PMP & P & PMP \\
\midrule
Dense ResNet-18, S0 only & \multicolumn{2}{c}{$87.2_{\pm2.5}$} & \multicolumn{2}{c}{$64.7_{\pm4.9}$} & \multicolumn{2}{c}{$91.4_{\pm5.6}$} & \multicolumn{2}{c}{$55.4_{\pm1.3}$} \\
Dense ResNet-18, S0+R+M+R & \multicolumn{2}{c}{$87.0_{\pm2.5}$} & \multicolumn{2}{c}{$64.2_{\pm4.7}$} & \multicolumn{2}{c}{$91.3_{\pm5.6}$} & \multicolumn{2}{c}{$55.5_{\pm1.1}$} \\
\midrule
Random & $80.1_{\pm3.6}$ & $81.1_{\pm3.0}$ & $57.8_{\pm6.7}$ & $58.4_{\pm4.6}$ & $85.6_{\pm8.1}$ & $86.8_{\pm7.7}$ & $46.6_{\pm1.3}$ & $46.7_{\pm1.7}$ \\
Filter $\ell_1$ & $81.0_{\pm3.1}$$^{\uparrow}$ & $81.5_{\pm3.0}$ & $57.9_{\pm6.2}$ & $59.4_{\pm5.7}$ & $85.8_{\pm8.3}$ & $86.4_{\pm7.7}$ & $48.7_{\pm0.9}$$^{\uparrow}$ & $48.1_{\pm0.8}$ \\
BN scale $|\gamma|$ & -- & $80.0_{\pm3.1}$$^{\downarrow}$ & -- & $56.8_{\pm4.9}$ & $85.1_{\pm8.9}$ & $85.1_{\pm8.7}$ & $44.4_{\pm1.7}$$^{\downarrow}$ & $45.5_{\pm1.0}$ \\
Taylor & $82.5_{\pm2.9}$$^{\uparrow}$ & $82.4_{\pm3.0}$$^{\uparrow}$ & $60.2_{\pm6.0}$$^{\uparrow}$ & $59.7_{\pm6.3}$ & $88.3_{\pm7.1}$$^{\uparrow}$ & $87.1_{\pm7.1}$ & $47.9_{\pm1.2}$ & $48.5_{\pm1.0}$$^{\uparrow}$ \\
Gradient $\mathcal{G}$ & -- & $81.4_{\pm3.3}$ & -- & $58.9_{\pm5.6}$ & $87.3_{\pm7.2}$ & $86.7_{\pm7.9}$ & $47.6_{\pm1.5}$ & $46.9_{\pm2.5}$ \\
Variance $\mathcal{V}$ & -- & $81.5_{\pm3.0}$ & -- & $57.3_{\pm4.4}$ & $86.5_{\pm6.9}$ & $86.5_{\pm7.3}$ & $47.2_{\pm1.4}$ & $47.3_{\pm1.3}$ \\
Fisher $\mathcal{D}$ & $79.6_{\pm3.1}$ & $79.9_{\pm2.9}$$^{\downarrow}$ & $58.1_{\pm6.5}$ & $58.4_{\pm6.0}$ & $85.2_{\pm8.9}$ & $85.3_{\pm8.0}$ & $46.8_{\pm0.8}$ & $46.4_{\pm1.3}$ \\
DACIS & $80.1_{\pm3.0}$ & $81.0_{\pm3.3}$ & $57.8_{\pm5.3}$ & $59.3_{\pm4.9}$ & $86.4_{\pm7.7}$ & $86.6_{\pm7.5}$ & $47.3_{\pm1.4}$ & $47.5_{\pm1.6}$ \\
\bottomrule
\end{tabular}
  \caption{Eight channel scores at $30\%$ of the parameters, 5-shot accuracy (\%) on novel classes. P prunes once and recovers; PMP is the prune, meta-train, prune schedule of Fig.~\ref{fig:overview}. Cells are the mean over \num{resisc45:seeds} class partitions with the standard deviation over partitions as subscript. PlantDoc columns test the PlantVillage-trained networks on field photographs.}
  \label{tab:scores}
\end{table*}

\paragraph{What 30\% of the parameters costs.}
Partition decides the level. The dense network scores \num{resisc45:dense_ce:novel_5}\% on RESISC45 at 5 shots, \num{eurosat13:dense_ce:novel_5}\% on EuroSAT-MS and \num{plantvillage:dense_ce:novel_5}\% on PlantVillage, with single partitions on EuroSAT-MS ranging from \num{eurosat13:dense_ce:novel_5:min}\% to \num{eurosat13:dense_ce:novel_5:max}\% (Table~\ref{tab:scores}). Keeping $30\%$ of the parameters with DACIS under PMP costs \num{d:rs:novel_5:dense-dacisPMP:abs} points on RESISC45 \ci{rs:novel_5:dense-dacisPMP}, \num{d:es:novel_5:dense-dacisPMP:abs} on EuroSAT-MS \ci{es:novel_5:dense-dacisPMP} and \num{d:pv:novel_5:dense-dacisPMP:abs} on PlantVillage \ci{pv:novel_5:dense-dacisPMP}; the brackets are $95\%$ intervals over partitions. On PlantDoc field photographs the same PlantVillage networks fall from \num{plantvillage:dense_ce:plantdoc_5}\% to \num{plantvillage:A_dacis_PMP_30:plantdoc_5}\%.

\paragraph{The choice among scores is a small effect, and DACIS is not ahead.}
Under PMP the best and the worst of the eight scores are \num{spread:rs:novel_5:PMP} points apart on RESISC45, \num{spread:es:novel_5:PMP} on EuroSAT-MS, \num{spread:pv:novel_5:PMP} on PlantVillage and \num{spread:pd:plantdoc_5:PMP} on PlantDoc. Random selection ranks \num{randrank:rs:novel_5:PMP}, \num{randrank:es:novel_5:PMP}, \num{randrank:pv:novel_5:PMP} and \num{randrank:pd:plantdoc_5:PMP} of eight (1 is best). Of the \num{cmp:total:PMP} paired comparisons of a computed score with random selection (seven scores, four settings), \num{cmp:above:PMP} have a $95\%$ interval above zero (\num{cmp:abovelist:PMP}) and \num{cmp:below:PMP} below it (\num{cmp:belowlist:PMP}); with that many intervals about \num{cmp:chance:PMP} would exclude zero by chance alone; after a Bonferroni correction the smallest adjusted $p$ value is \num{cmp:bonfabove:PMP} among the comparisons that favour a computed score and \num{cmp:bonf:PMP} among all of them, so the deficits of the worst scores (BatchNorm scale and Fisher ratio on RESISC45) are the only differences that survive it. The largest mean advantage is \num{cmp:maxgain:PMP} points, and the intervals have a half-width of \num{cmp:halfw:PMP} points on average (up to \num{cmp:halfwmax:PMP}), so five partitions cannot rule out differences of that size either way. Under PMP the first-order Taylor term is the only score ahead of random selection, on two settings: by \num{d:rs:novel_5:taylor-random:PMP} points on RESISC45 \ci{rs:novel_5:taylor-random:PMP} and \num{d:pd:plantdoc_5:taylor-random:PMP} on PlantDoc \ci{pd:plantdoc_5:taylor-random:PMP} under PMP, and by \num{d:rs:novel_5:taylor-random:P} on RESISC45 under P. DACIS minus random is \num{d:rs:novel_5:dacis-random:PMP} points on RESISC45 \ci{rs:novel_5:dacis-random:PMP}, \num{d:es:novel_5:dacis-random:PMP} on EuroSAT-MS \ci{es:novel_5:dacis-random:PMP} and \num{d:pv:novel_5:dacis-random:PMP} on PlantVillage \ci{pv:novel_5:dacis-random:PMP}, and DACIS minus Taylor is \num{d:rs:novel_5:dacis-taylor:PMP} on RESISC45 \ci{rs:novel_5:dacis-taylor:PMP} (Table~\ref{tab:pairs}). On the PlantVillage and EuroSAT-MS partitions the ranking of the scores hardly repeats (Kendall's $W$ \num{kw:pv:novel_5:PMP} and \num{kw:es:novel_5:PMP}, against about $1/\num{resisc45:seeds}$ for unrelated rankings), while on RESISC45 it does (\num{kw:rs:novel_5:PMP}) because Taylor is usually first and Fisher or BatchNorm scale last (Fig.~\ref{fig:seeds}). The channels the scores keep are far from identical (Jaccard overlap in the supplement), so the small differences do not come from the scores choosing the same channels.

\begin{figure*}[t]
  \centering
  \includegraphics[width=\textwidth]{generated/fig_seeds.pdf}
  \caption{Accuracy of each score at $30\%$ of the parameters (schedule PMP) on each class partition, relative to the mean over scores.}
  \label{fig:seeds}
\end{figure*}

\begin{table*}[t]
  \centering\small\setlength{\tabcolsep}{4pt}
  \begin{tabular}{@{}lcccc@{}}
\toprule
Difference (a $-$ b), points & RESISC45 & EuroSAT-MS & PlantVillage & PlantDoc \\
\midrule
\multicolumn{5}{@{}l}{\emph{Scores, schedule P}} \\
DACIS $-$ Random & $\ensuremath{-}0.1$ {\scriptsize[\ensuremath{-}0.9,+0.8]} & $+0.1$ {\scriptsize[\ensuremath{-}2.6,+2.8]} & $+0.8$ {\scriptsize[\ensuremath{-}1.3,+3.0]} & $+0.7$ {\scriptsize[\ensuremath{-}1.2,+2.6]} \\
DACIS $-$ Filter $\ell_1$ & $\ensuremath{-}0.9$ {\scriptsize[\ensuremath{-}1.5,\ensuremath{-}0.3]} & $\ensuremath{-}0.1$ {\scriptsize[\ensuremath{-}2.2,+2.1]} & $+0.6$ {\scriptsize[\ensuremath{-}1.9,+3.1]} & $\ensuremath{-}1.4$ {\scriptsize[\ensuremath{-}3.3,+0.5]} \\
DACIS $-$ Taylor & $\ensuremath{-}2.4$ {\scriptsize[\ensuremath{-}3.0,\ensuremath{-}1.7]} & $\ensuremath{-}2.4$ {\scriptsize[\ensuremath{-}4.9,+0.1]} & $\ensuremath{-}1.9$ {\scriptsize[\ensuremath{-}2.9,\ensuremath{-}0.9]} & $\ensuremath{-}0.6$ {\scriptsize[\ensuremath{-}1.6,+0.4]} \\
\multicolumn{5}{@{}l}{\emph{Scores, schedule PMP}} \\
DACIS $-$ Random & $\ensuremath{-}0.0$ {\scriptsize[\ensuremath{-}0.8,+0.7]} & $+0.9$ {\scriptsize[\ensuremath{-}1.3,+3.0]} & $\ensuremath{-}0.3$ {\scriptsize[\ensuremath{-}2.3,+1.8]} & $+0.8$ {\scriptsize[\ensuremath{-}0.4,+2.1]} \\
DACIS $-$ Filter $\ell_1$ & $\ensuremath{-}0.4$ {\scriptsize[\ensuremath{-}1.0,+0.1]} & $\ensuremath{-}0.1$ {\scriptsize[\ensuremath{-}1.6,+1.4]} & $+0.2$ {\scriptsize[\ensuremath{-}1.3,+1.7]} & $\ensuremath{-}0.7$ {\scriptsize[\ensuremath{-}2.4,+1.1]} \\
DACIS $-$ Taylor & $\ensuremath{-}1.4$ {\scriptsize[\ensuremath{-}1.9,\ensuremath{-}0.9]} & $\ensuremath{-}0.4$ {\scriptsize[\ensuremath{-}2.7,+1.9]} & $\ensuremath{-}0.6$ {\scriptsize[\ensuremath{-}1.9,+0.8]} & $\ensuremath{-}1.1$ {\scriptsize[\ensuremath{-}2.4,+0.3]} \\
\multicolumn{5}{@{}l}{\emph{Episodic stage on the dense network}} \\
Dense+M $-$ Dense & $\ensuremath{-}0.3$ {\scriptsize[\ensuremath{-}0.4,\ensuremath{-}0.1]} & $\ensuremath{-}0.8$ {\scriptsize[\ensuremath{-}1.5,+0.0]} & $\ensuremath{-}0.3$ {\scriptsize[\ensuremath{-}0.6,+0.1]} & $\ensuremath{-}0.0$ {\scriptsize[\ensuremath{-}0.4,+0.3]} \\
\bottomrule
\end{tabular}
  \caption{Paired differences in points (mean over partitions, with the $95\%$ $t$-interval over partitions in brackets); (a $-$ b) is positive when a is more accurate. Test episodes are shared inside a partition. A dash marks a comparison that was not run on that dataset.}
  \label{tab:pairs}
\end{table*}

\paragraph{Episodic training did not help the dense network, and the pruning schedules gain at most a point.}
Adding $1000$ episodes ($800$ on EuroSAT-MS) of prototypical training after cross-entropy pre-training changes 5-shot accuracy of the dense network by \num{d:rs:novel_5:densem-dense} points on RESISC45 \ci{rs:novel_5:densem-dense}, \num{d:es:novel_5:densem-dense} on EuroSAT-MS \ci{es:novel_5:densem-dense} and \num{d:pv:novel_5:densem-dense} on PlantVillage \ci{pv:novel_5:densem-dense}, at the learning rate chosen on validation classes (\num{tune:best}). In the validation sweep (supplement) 5-shot accuracy broadly falls as the learning rate grows, on both sets that have validation classes, and no value beats stopping after S0; our conclusion is for this recipe, with frozen BatchNorm and at most the smallest learning rate of the grid. Inside pruning schedules (Table~\ref{tab:pairs}) PMP is ahead of P by \num{d:rs:novel_5:PMP-P} points on RESISC45 \ci{rs:novel_5:PMP-P} and level elsewhere. The gradient probe P$\to$G$\to$P, which collects $G_{\mathrm{meta}}$ without updating a weight in the meta stage, is ahead by \num{d:rs:novel_5:PGP-P} \ci{rs:novel_5:PGP-P}, as much as PMP. The gain therefore comes from the extra pruning step and recovery epoch that both schedules contain, not from episodic learning. Guo et al.~\cite{guo2020bscd} found that meta-learning methods trail simple fine-tuning on their crop-disease and EuroSAT targets.

\begin{table*}[t]
  \centering\small\setlength{\tabcolsep}{4pt}
  \begin{tabular}{@{}lrrrrrrr@{}}
\toprule
Model & Params (M) & MACs (G) & CPU ms & GPU img/s & RESISC45 & EuroSAT-MS & PlantVillage \\
\midrule
ResNet-18, dense & 11.18 & 0.592 & 10.8 & 36520 & $86.9_{\pm3.4}$ & $62.1_{\pm5.3}$ & $91.4_{\pm5.6}$ \\
ResNet-18, 50\% (DACIS) & 5.59 & 0.398 & 8.6 & 38034 & $83.8_{\pm3.7}$ & $58.6_{\pm5.7}$ & $88.6_{\pm7.2}$ \\
ResNet-18, 30\% (DACIS) & 3.35 & 0.251 & 6.6 & 40431 & $80.5_{\pm4.2}$ & $57.1_{\pm4.9}$ & $86.6_{\pm7.5}$ \\
ResNet-18, 20\% (DACIS) & 2.24 & 0.178 & 5.7 & 41645 & $77.4_{\pm4.9}$ & $54.5_{\pm8.1}$ & $83.8_{\pm8.2}$ \\
ResNet-18, 10\% (DACIS) & 1.11 & 0.105 & 4.6 & 45865 & $71.9_{\pm4.8}$ & $51.4_{\pm7.5}$ & $78.2_{\pm10.2}$ \\
\midrule
MobileNetV3-Small, dense & 0.93 & 0.018 & 2.3 & 63688 & $88.1_{\pm3.3}$ & $72.4_{\pm4.9}$ & $91.1_{\pm5.9}$ \\
MobileNetV2, dense & 2.22 & 0.098 & 4.3 & 24416 & $86.9_{\pm3.1}$ & $65.4_{\pm3.5}$ & $91.2_{\pm6.3}$ \\
EfficientNet-B0, dense & 4.01 & 0.126 & 6.3 & 18928 & $90.5_{\pm2.4}$ & $69.9_{\pm5.1}$ & $93.8_{\pm5.2}$ \\
\bottomrule
\end{tabular}
  \caption{Cost and 5-shot accuracy (\%). The pruned architecture is the same for every score, so one cost row serves all scores; accuracy is for DACIS under PMP. Cost is for a $128{\times}128$ RGB input: CPU is single-thread latency at batch size 1 on \num{hw-cpu}, GPU throughput is at batch size 128 on \num{hw-gpu}. Latency is measured, not a statement about embedded devices. The accuracy columns use each dataset's own input size and bands. The compact networks are trained with the same recipe and are not pruned.}
  \label{tab:budget}
\end{table*}

\begin{figure*}[t]
  \centering
  \includegraphics[width=\textwidth]{generated/fig_tradeoff.pdf}
  \caption{5-shot accuracy against multiply-accumulates (log scale), mean over \num{resisc45:seeds} partitions. The orange line connects the ResNet-18 pruned with DACIS to $50$, $30$, $20$ and $10\%$ of the parameters; the compact networks are not pruned. Accuracy against budget for four scores is in the supplement.}
  \label{fig:budget}
\end{figure*}

\paragraph{Budget matters more than score, and a small pretrained network beats all of them.}
Figure~\ref{fig:budget} shows accuracy falling steadily with the budget and the compact networks above and to the left of every pruned ResNet-18: on RESISC45 DACIS drops from \num{resisc45:C_dacis_PMP_50:novel_5}\% at $50\%$ of the parameters to \num{resisc45:C_dacis_PMP_10:novel_5}\% at $10\%$, a larger step than the whole spread among scores. Pruning saves less time than operations (Table~\ref{tab:budget}): at $30\%$ of the parameters a ResNet-18 needs \num{resisc45:A_l1_PMP_30:macs_128} G multiply-accumulates at $128{\times}128$ against \num{resisc45:dense_ce:macs_128} G for the dense network, which is \num{cost:macs-red:30}\% fewer, but its single-thread CPU latency is only \num{cost:cpu-red:30}\% lower. We did not profile layer by layer. The stem, the shortcuts and the residual widths are untouched, and at batch size one the small layers may be limited by memory traffic and call overhead more than by arithmetic; either would explain the gap, and we cannot tell which. The off-the-shelf compact networks are a strong baseline that pruning studies in this setting should include. MobileNetV3-Small, with \num{cost:mobilenet_v3_small:params} M parameters, \num{cost:mobilenet_v3_small:macs} G multiply-accumulates and \num{cost:mobilenet_v3_small:cpu} ms on the CPU, scores \num{resisc45:K_mobilenet_v3_small:novel_5}\% on RESISC45, \num{eurosat13:K_mobilenet_v3_small:novel_5}\% on EuroSAT-MS and \num{plantvillage:K_mobilenet_v3_small:novel_5}\% on PlantVillage, which is \num{d:rs:novel_5:mnv3-dacisPMP} points above the $30\%$ ResNet-18 on RESISC45 \ci{rs:novel_5:mnv3-dacisPMP} and \num{d:es:novel_5:mnv3-dacisPMP} on EuroSAT-MS \ci{es:novel_5:mnv3-dacisPMP}, and above the unpruned ResNet-18 on EuroSAT-MS (\num{eurosat13:dense_ce:novel_5}\%). EfficientNet-B0 is higher still on RESISC45 and PlantVillage.

\paragraph{Sentinel-2 bands.}
On EuroSAT-MS the dense network with all thirteen bands scores \num{eurosat13:dense_ce:novel_5}\% at 5 shots against \num{eurosat3:dense_ce:novel_5}\% for the same network on the three visible bands (Table~\ref{tab:eurosat}). In this naive adaptation, with five base classes and ImageNet filters repeated over the extra bands, the additional bands lowered novel-class accuracy; the interval over partitions is wide, and we do not conclude that spectral bands are unhelpful for few-shot EuroSAT. DACIS minus filter $\ell_1$ under PMP is \num{d:es:novel_5:dacis-l1:PMP} points with thirteen bands and \num{d:e3:novel_5:dacis-l1:PMP} with three.

\begin{table}[t]
  \centering\scriptsize\setlength{\tabcolsep}{2pt}
  \begin{tabular}{@{}lrrr@{}}
\toprule
EuroSAT-MS, novel classes & 1-shot & 5-shot & 10-shot \\
\midrule
13 bands, dense (S0) & $49.4_{\pm5.5}$ & $62.1_{\pm5.3}$ & $64.9_{\pm5.0}$ \\
RGB bands, dense (S0) & $53.7_{\pm4.3}$ & $68.7_{\pm3.9}$ & $71.4_{\pm3.6}$ \\
13 bands, dense (S0+R+M+R) & $49.0_{\pm5.2}$ & $61.7_{\pm5.1}$ & $64.3_{\pm5.0}$ \\
RGB bands, dense (S0+R+M+R) & $53.1_{\pm4.5}$ & $68.1_{\pm3.8}$ & $70.9_{\pm3.4}$ \\
13 bands, filter $\ell_1$, PMP 30\% & $46.1_{\pm6.5}$ & $57.0_{\pm6.0}$ & $59.3_{\pm5.2}$ \\
RGB bands, filter $\ell_1$, PMP 30\% & $49.8_{\pm2.4}$ & $63.9_{\pm2.6}$ & $66.6_{\pm2.3}$ \\
13 bands, DACIS, PMP 30\% & $46.0_{\pm5.4}$ & $57.1_{\pm4.9}$ & $59.4_{\pm4.1}$ \\
RGB bands, DACIS, PMP 30\% & $50.2_{\pm3.5}$ & $63.6_{\pm3.9}$ & $66.2_{\pm3.7}$ \\
13 bands, MobileNetV3-Small & $57.7_{\pm7.7}$ & $72.4_{\pm4.9}$ & $75.7_{\pm4.5}$ \\
RGB bands, MobileNetV3-Small & $58.3_{\pm3.7}$ & $74.7_{\pm4.1}$ & $78.0_{\pm4.1}$ \\
\midrule
13 bands $-$ RGB, dense (5-shot) & \multicolumn{3}{c}{$\ensuremath{-}6.6$ {\scriptsize[\ensuremath{-}17.4,+4.2]}} \\
13 bands $-$ RGB, DACIS PMP (5-shot) & \multicolumn{3}{c}{$\ensuremath{-}6.4$ {\scriptsize[\ensuremath{-}15.5,+2.6]}} \\
\bottomrule
\end{tabular}
  \caption{EuroSAT-MS, novel classes: thirteen bands against the three visible bands. The extra bands lower accuracy in this adaptation.}
  \label{tab:eurosat}
\end{table}

\paragraph{Recovery and the parts of DACIS (PlantVillage).}
A longer recovery after pruning did little: with six epochs instead of one, random selection and filter $\ell_1$ reach \num{plantvillage:R_random_P_30_r6:novel_5}\% and \num{plantvillage:R_l1_P_30_r6:novel_5}\% under schedule P on PlantVillage, against \num{plantvillage:P_random:novel_5}\% and \num{plantvillage:P_l1:novel_5}\% with one. Of the \num{abl:count} variants of DACIS in the supplement (single components, pairs, no meta-gradient term, other depth profile or scaling), the largest mean change in PlantVillage PMP accuracy is \num{abl:max} points, \num{abl:maxsign} the full score, for \num{abl:maxname}; \num{abl:nsig} of them have a $95\%$ interval clear of zero. \label{sec:lambda}On the validation classes, the eight weightings of Eq.~\eqref{eq:dacis} that we tried span \num{lam:val5:worst}\% to \num{lam:val5:best}\% at 5 shots against \num{lam:val5:default}\% for the weights fixed in advance, so we left $\lambda$ alone; this is a sensitivity check on validation classes, not a tuning protocol.
